\subsection{Products of weak topologies}

PROPOSITION 8. --- Let \((\mathbf{F}_{\iota}, \mathbf{G}_{\iota})_{\iota \in \mathbf{I}}\) be a family of pairs of spaces in duality. Let \(\mathbf{F} = \prod_{\iota \in \mathbf{I}} \mathbf{F}_{\iota}\) be the product space of the \(\mathbf{F}_{\iota}\) and \(\mathbf{G} = \bigoplus_{\iota \in \mathbf{I}} \mathbf{G}_{\iota}\) be the direct sum of the \(\mathbf{G}_{\iota}\) . If, for all \(x = (x_{\iota}) \in \mathbf{F}\) and all \(y = (y_{\iota}) \in \mathbf{G}\) , we write \(\langle x, y \rangle = \sum_{\iota \in \mathbf{I}} \langle x_{\iota}, y_{\iota} \rangle\) (a sum which has only finitely many non-zero terms) then the topology \(\sigma(\mathbf{F}, \mathbf{G})\) (relative to the bilinear form \((x, y) \mapsto \langle x, y \rangle\) ) is the product of the topologies \(\sigma(\mathbf{F}_{\iota}, \mathbf{G}_{\iota})\) .

For, given a topology \(\mathcal{T}\) on \(\mathbf{F}\) ; in order that, for all \(y \in \mathbf{G}\) , the linear form \(x \mapsto \langle x, y \rangle\) should be continuous for \(\mathcal{T}\) , it is necessary and sufficient, by the definition of \(\langle x, y \rangle\) , that each of the mappings \(x \mapsto \langle \mathrm{pr}_i x, y_i \rangle\) should be continuous for \(\mathcal{T}\) , where \(i\) is arbitrary in I and \(y_i\) in \(\mathbf{G}_i\) ; but this means that each of the mappings \(\mathrm{pr}_i\) of F in \(\mathbf{F}_i\) is continuous for \(\mathcal{T}\) and for \(\sigma(\mathbf{F}_i, \mathbf{G}_i)\) (I, p.~10, cor. 1); this completes the demonstration.

Remark. --- The duality between F and G is separating in F (resp. in G) if and only if for all \(\iota\in I\) , the duality between \(F_{\iota}\) and \(G_{\iota}\) is separating in \(F_{\iota}\) (resp. in \(G_{\iota}\) ). If the duality between F and G is separating in F (resp. G), then, in F (resp. G), the subspace orthogonal to one \(G_{\iota}\) (resp. \(F_{\iota}\) ), canonically identified with a subspace of G (resp. F) is the subspace of the product of the \(F_{\kappa}\) where \(\kappa\neq\iota\) (resp. the direct sum of the \(G_{\kappa}\) such that \(\kappa\neq\iota\) ).

COROLLARY 1. --- Let F and G be two vector spaces in separating duality. If the space F (with \(\sigma(F, G)\) ) is the direct topological sum of two subspaces M, N then the space G (with \(\sigma(G, F)\) ) is the direct topological sum of the subspaces M°, N° orthogonal respectively to M and N.

For let \(p: \mathbf{F} \to \mathbf{M}\) , \(q: \mathbf{F} \to \mathbf{N}\) be the projectors corresponding to the decomposition of \(\mathbf{F}\) into the direct sum of \(\mathbf{M}\) and \(\mathbf{N}\) ; in these conditions the mapping \((p, q): \mathbf{F} \to \mathbf{M} \times \mathbf{N}\) is a topological isomorphism. If \(\mathbf{M}_1 = \mathbf{G} / \mathbf{M}^\circ\) , \(\mathbf{N}_1 = \mathbf{G} / \mathbf{N}^\circ\) , then the topologies on \(\mathbf{M}\) and \(\mathbf{N}\) (induced by that of \(\mathbf{F}\) ) are identical with \(\sigma(\mathbf{M}, \mathbf{M}_1)\) , \(\sigma(\mathbf{N}, \mathbf{N}_1)\) respectively (II, p.~48, prop. 7). The mapping \(^t(p, q): \mathbf{M}_1 \times \mathbf{N}_1 \to \mathbf{G}\) is a topological isomorphism when we give \(\mathbf{M}_1, \mathbf{N}_1\) and \(\mathbf{G}\) the topologies \(\sigma(\mathbf{M}_1, \mathbf{M})\) , \(\sigma(\mathbf{N}_1, \mathbf{N})\) and \(\sigma(\mathbf{G}, \mathbf{F})\) , by prop. 8. Under this mapping \(\mathbf{M}_1\) (resp. \(\mathbf{N}_1\) ) has as its image in \(\mathbf{G}\) the subspace \(\mathbf{N}^\circ\) (resp. \(\mathbf{M}^\circ\) ), and the topology \(\sigma(\mathbf{M}_1, \mathbf{M})\) (resp. \(\sigma(\mathbf{N}_1, \mathbf{N})\) ) has as its image the topology induced on \(\mathbf{N}^\circ\) (resp. \(\mathbf{M}^\circ\) ) by \(\sigma(\mathbf{G}, \mathbf{F})\) , from which the corollary follows.

COROLLARY 2. --- Let \((e_{\iota})_{\iota \in \mathbb{J}}\) be a basis of the vector space F with dual F*, and let \(u: \mathbf{R}^{(I)} \to \mathbf{F}\) be an (algebraic) isomorphism defined by this basis. Then the transposed mapping \({}^t u: \mathbf{F}^* \to \mathbf{R}^I\) is a topological isomorphism when F* carries the topology \(\sigma(\mathbf{F}^*, \mathbf{F})\) and \(\mathbf{R}^I\) the product topology.

We know (A, II, § 2.6, prop. 10) that \({}^t u\) is a bijection, and that if for a \(x^* \in \mathbf{F}^*\) , we put \(\langle e_i, x^* \rangle = \xi_i^*\) for all \(i \in I\) , then the image \({}^t u(x^*)\) is the vector \((\xi_i^*)\) of \(\mathbf{R}^I\) , so that, for all \(x = \sum_{i} \xi_i e_i\) in \(F\) , we have \(\langle x, x^* \rangle = \sum_{i \in I} \xi_i \xi_i^*\) . The corollary then follows from this formula and prop. 8.

